% Solutions/hints for ch21 exercises (assembled into Appendix C)
\begin{solution}{exr:ch21-condensed}
By induction on $k$: $\state_1=\mat{A}\state_0+\mat{B}\ctrl_0$, and if $\state_k=\mat{A}^k\state_0+\sum_{j=0}^{k-1}\mat{A}^{k-1-j}\mat{B}\ctrl_j$ then $\state_{k+1}=\mat{A}\state_k+\mat{B}\ctrl_k=\mat{A}^{k+1}\state_0+\sum_{j=0}^{k}\mat{A}^{k-j}\mat{B}\ctrl_j$. Stacking $k=1,\dots,N$ gives $\mat{X}=\mat{S}_x\state_0+\mat{S}_u\mat{U}$ with the $k$-th block row of $\mat{S}_x$ equal to $\mat{A}^k$ and the block $(k,j)$ of $\mat{S}_u$ equal to $\mat{A}^{k-1-j}\mat{B}$ for $j<k$ and $\mat{0}$ otherwise. Substituting into the cost, $(\mat{X}-\mat{X}^{\mathrm{ref}})\T\bar{\mat{Q}}(\mat{X}-\mat{X}^{\mathrm{ref}})+\mat{U}\T\bar{\mat{R}}\mat{U}=\mat{U}\T(\mat{S}_u\T\bar{\mat{Q}}\mat{S}_u+\bar{\mat{R}})\mat{U}+2(\mat{S}_x\state_0-\mat{X}^{\mathrm{ref}})\T\bar{\mat{Q}}\mat{S}_u\mat{U}+\text{const}$, so $\mat{H}=2(\mat{S}_u\T\bar{\mat{Q}}\mat{S}_u+\bar{\mat{R}})$ and $\vect{f}=2\mat{S}_u\T\bar{\mat{Q}}(\mat{S}_x\state_0-\mat{X}^{\mathrm{ref}})$. $\mat{H}$ is positive definite because $\bar{\mat{R}}$ is and $\mat{S}_u\T\bar{\mat{Q}}\mat{S}_u$ is positive semidefinite; hence the QP is strictly convex and has a unique minimiser. The velocity bound $-v_{\max}\le v_{k,i}\le v_{\max}$ becomes $-v_{\max}-(\mat{S}_x\state_0)_{v_{k,i}}\le(\mat{S}_u)_{v_{k,i}}\mat{U}\le v_{\max}-(\mat{S}_x\state_0)_{v_{k,i}}$, one row of $\mat{S}_u$ per velocity component and step.
\end{solution}
\begin{solution}{exr:ch21-halfplane}
(a) If $\vect{n}\T(\pos-\vect{o})\ge r$ with $\norm{\vect{n}}=1$, the Cauchy--Schwarz inequality gives $\norm{\pos-\vect{o}}\ge\vect{n}\T(\pos-\vect{o})\ge r$, so every point of the half-plane is outside the disc; the half-plane is the tangent half-plane at the point $\vect{o}+r\vect{n}$ and it also excludes the collision-free points ``behind'' the disc as seen from the linearisation point. (b) The half-planes $\vect{n}_m\T(\pos-\vect{c})\le e\cos(\pi/M)$ for $M$ equally spaced unit directions describe the regular $M$-gon inscribed in the disc of radius $e$: its apothem is $e\cos(\pi/M)$ and its vertices lie on the circle, so the polygon is inside the disc and the constraint is conservative. The largest loss is at the flat sides, $e(1-\cos(\pi/M))$: $7.6\,\%$ of $e$ for $M=8$ and $1.9\,\%$ for $M=16$. (c) In three dimensions replace the disc by a ball and the polygon by a polyhedron, e.g.\ the $26$ directions $(\pm1,\pm1,\pm1)$, $(\pm1,\pm1,0)$ and permutations, normalised; the apothem is then the smallest distance from the centre to a face and must be computed once.
\end{solution}
\begin{solution}{exr:ch21-horizon}
A drone at speed $v$ needs the time $v/a_{\max}$ and the distance $v^2/(2a_{\max})$ to stop (\cref{ch:ch02}). If the horizon $N\dt$ is shorter than the stopping time, the QP is asked to guarantee separation at steps it cannot influence enough, and a hard constraint that first appears at the end of the horizon can already be impossible to satisfy. With $v_{\max}=2$~m/s and $a_{\max}=2$~m/s$^2$ the stopping time is $1$~s, i.e.\ $N\ge10$ at $\dt=0.1$~s; at the cruising speed $1$~m/s of the example it is $0.5$~s, $N\ge5$, which is why the generated experiment shows failures at $N=3$ and none from $N=4$ on. Add the reaction latency and the prediction horizon of the intruder to this bound in practice.
\end{solution}
\begin{solution}{exr:ch21-slack}
Write the soft problem as $\min_{\mat{U},\vect{s}}J(\mat{U})+w\sum_k s_k$ subject to $g_k(\mat{U})\le s_k$, $s_k\ge0$. Its Lagrangian coincides with that of the hard problem plus the terms $(w-\lambda_k-\mu_k)s_k$, where $\lambda_k\ge0$ is the multiplier of the avoidance row and $\mu_k\ge0$ that of $s_k\ge0$. At a solution of the hard problem with multipliers $\lambda_k^*$, choosing $s_k=0$ and $\mu_k=w-\lambda_k^*\ge0$ satisfies all optimality conditions of the soft problem as soon as $w\ge\max_k\lambda_k^*$; because the problem is convex these conditions are sufficient, so the soft solution equals the hard one. If $w<\lambda_k^*$ for some $k$, the penalty is cheaper than the tracking cost it saves and the solver buys a violation $s_k>0$. In the worked example the largest multiplier is about $16$, so $w=50$ and $w=100$ reproduce the hard trajectory exactly, $w=10$ is short of it by $2$~mm of separation, and $w=1$ lets the drone cut $0.22$~m into the safety disc.
\end{solution}
